% Finite normalization descent from the independently audited generic-surface theorem.
\begin{corollary}[Proved: finite punctures in arbitrary-dimensional proper schemes]
\label{D:nonnormal-finite-puncture-quantitative-transfer}
Let \(S\) be a proper scheme over an algebraically closed field
of characteristic \(p>0\).  Let \(T\) be a finite closed subset
which has codimension at least two in every irreducible component
it meets, and let \(U\) be a connected component of \(S\setminus T\).
Write \(Q=\pi_1(U)\).  Then the following statements hold.
\begin{enumerate}
\item There is a constant \(c_{S,T}\), independent of the finite
coefficient field \(\Lambda\), such that
\[
 \dim_\Lambda H^1(Q,M)\leq c_{S,T}\dim_\Lambda M
\]
for every finite-dimensional continuous \(\Lambda[Q]\)-module \(M\).
The corresponding bound for an open subgroup \(J\) is
\(c_{S,T}[Q:J]\dim_\Lambda M\).
In particular, the maximal prosolvable quotient of \(Q\) is
topologically finitely generated.
\item There is a constant \(C_{S,T}\) such that, for every \(q=p^f\)
and every \(r\geq1\),
\[
 \#\left(\operatorname{Hom}_{\mathrm{cts}}(Q,\operatorname{GL}_r(\mathbf F_q))
          /\operatorname{GL}_r(\mathbf F_q)\right)
 \leq q^{C_{S,T}r^2}.
 \tag{FP1}
\]
\item Fix \(B\geq1\).  For a finite subgroup
\(\widetilde G\subset\operatorname{GL}_r(\mathbf F_q)\), suppose
its scalar subgroup has order at most \(B\), and denote its image
in \(\operatorname{PGL}_r(\mathbf F_q)\) by \(G\).
There is a constant \(C_{S,T,B}\), independent of \(q,r,\widetilde G\),
such that the number of homomorphisms \(Q\to G\), up to conjugacy
in the ambient \(\operatorname{PGL}_r(\mathbf F_q)\), is at most
\[
                         q^{C_{S,T,B}r^2}.                \tag{FP2}
\]
Consequently, for any fixed \(\epsilon>0\), the quotient formed
from all distinct nonabelian simple quotients in such a family
with \(|G|\geq q^{\epsilon r^2}\) is topologically finitely generated.
\end{enumerate}
In particular the bounded-centre linear and unitary families,
and the natural symplectic and orthogonal families established
earlier, satisfy these assertions.  No degree-two cohomology
bound for \(Q\), nor full finite generation of \(Q\), is asserted.
\end{corollary}

\begin{proof}
Reduction does not affect finite etale covers or their local
systems.  Normalize \(S_{\mathrm{red}}\), and let \(V\) be the
inverse image of \(U\).  Excellence of schemes of finite type
over a field makes \(V\to U\) finite and surjective.  Write its
finitely many connected components as \(V_1,\ldots,V_t\).
Each \(V_i\) is a normal proper integral component with finitely
many points removed.  Indeed that complement is the inverse
image of the finite set \(T\); connectedness of a nonempty open
of an integral scheme ensures that a normalized component lies
over a single connected component of \(S\setminus T\).
The codimension hypothesis makes the components of dimension
at most one proper.

For a normal component of dimension at least two,
Theorem~\ref{D:generic-surface-finite-punctures} gives a continuous surjection
\[
 \pi_1(C_i\setminus F_i)\longrightarrow\pi_1(V_i),
 \tag{FP3}
\]
where \(C_i\) is a normal proper surface over an algebraically
closed extension of the ground field and \(F_i\) is finite.
The surface is fixed once \(V_i\) is fixed; it does not depend
on a finite covering group, a representation, or its coefficient
field.  Inflation along (FP3) is injective on \(H^1\).
Precomposition is also injective on representation classes,
whether ordinary linear or ambient projective conjugacy is
used: an equality after a surjective precomposition is an
equality on the target group.

Theorems~\ref{D:punctured-surface-all-characteristics-H1}
and~\ref{D:global-quadratic-representation-bound}, and
Proposition~\ref{D:bounded-scalar-projective-growth}, therefore
give all three bounds on each \(V_i\).  For components of
dimension at most one we use the same already proved
dimension-at-most-two statements with empty omitted set.
Thus there are constants \(c_i,C_i,C_{i,B}\) with the required
uniformity on every normalized component.

Put
\[
                   c=\#\pi_0(V\times_U V).
\]
This is finite, since the double product is noetherian.
We explain both the linear and nonabelian descent estimates;
neither uses cohomology in degree two.

\par\smallskip\noindent\textbf{The first-cohomology bound.}\quad
Regard \(M\) as a finite locally constant sheaf of
\(\Lambda\)-vector spaces on \(U\).  A class in the kernel of
\[
 H^1(U,M)\longrightarrow\bigoplus_i H^1(V_i,M|_{V_i})
\]
is an \(M\)-torsor with a chosen trivialization on \(V\).
The difference of its two pulled-back trivializations is a
section of \(M\) on \(V\times_U V\), satisfying the usual
cocycle condition.  Changing the trivialization adds a
coboundary.  Effective finite-proper descent of finite etale
schemes therefore identifies this kernel with a quotient of
a subspace of \(H^0(V\times_U V,M)\)
\citep[Expose~IX, Theorem~4.7]{C_SGA1}.
A section of a finite local system on a connected scheme is
determined by its value in one geometric stalk.  Consequently
\[
 \dim_\Lambda H^0(V\times_U V,M)\leq c\dim_\Lambda M,
\]
and the desired bound follows with
\(c_{S,T}=c+\sum_i c_i\).
Shapiro's lemma gives the open-subgroup statement.
The prosolvable conclusion is the abelian-crown consequence
of Proposition~\ref{D:cohomology-crown-criterion}.

\par\smallskip\noindent\textbf{Linear representation classes.}\quad
Set \(P=\operatorname{GL}_r(\mathbf F_q)\).
Its representation classes are the isomorphism classes of
torsors for the constant finite group \(P\).
There are at most \(q^{(\sum_i C_i)r^2}\) possible tuples of
their restriction classes to \(V_i\).
Fix representative torsors for one such tuple.
The Isom sheaf between their two pullbacks to \(V\times_U V\)
is finite etale, and each geometric fibre has \(|P|\) elements.
A section on a connected component is determined by one
geometric stalk.  Hence there are at most \(|P|^c\) possible
descent isomorphisms.  The cocycle condition only reduces this
number, and passing to isomorphism classes only reduces it
again.  Finite-proper descent, including full faithfulness for
the group action maps, is the same cited theorem.
Since \(|P|\leq q^{r^2}\), we obtain (FP1).

\par\smallskip\noindent\textbf{Ambient projective conjugacy.}\quad
Now put \(P=\operatorname{PGL}_r(\mathbf F_q)\).
A homomorphism \(Q\to G\) defines a \(P\)-torsor by extension
of structure group.  Two such homomorphisms give the same
torsor class exactly when they are conjugate in \(P\).
This is an equivalence relation on the \(G\)-valued
homomorphisms even when the ambient group does not normalize
the chosen subgroup \(G\).
Restrictions of these torsors to each \(V_i\) still arise
from \(G\)-valued homomorphisms, so their possible tuples
number at most \(q^{(\sum_i C_{i,B})r^2}\).
For each fixed tuple we may count all \(P\)-torsor descent
data, including those whose descended monodromy leaves \(G\);
this is a valid upper bound.  The preceding Isom argument
again contributes at most \(|P|^c\leq q^{cr^2}\), proving (FP2).

Distinct epimorphism kernels are separated by ambient
projective-conjugacy classes.  The order lower bound in the
statement converts (FP2) into a fixed polynomial bound on
the multiplicity of each simple quotient.
The generating-tuple argument already used in
Proposition~\ref{D:bounded-scalar-projective-growth} proves
finite generation of the stated simple-product quotient.
The argument is entirely in degree zero and degree one
descent; it does not infer an \(H^2\) bound from (FP3).
\end{proof}
